
\documentclass{article}
\usepackage{amsmath}
\title{FT1.1 — Mathematical Model for Sudoku / Generalized Sudoku}
\author{Sistema Operador Unificador Fenomenológico General}
\begin{document}
\maketitle

\section*{Equation}
\begin{equation}
S^* = \arg\min_S [ \Phi(S) + \lambda \sum_{x,y} |z_{k+1}-z_k| ]
\end{equation}

Solution condition:
\begin{equation}
\Phi(S) = 0
\end{equation}

\section*{Variables}
\begin{itemize}
\item $S(x,y,z)$ : board state
\item $S^*$ : optimal solution
\item $\Phi(S)$ : constraint violation function
\item $\lambda$ : depth regulation factor
\item $x,y$ : board coordinates
\item $z$ : logical exploration depth
\item $\beta \in \{-1,1\}$ : depth displacement bias
\end{itemize}

\section*{Domain}
$D = \{1,2,\dots,n^2\}$

Classic Sudoku:
\begin{itemize}
\item Board: $9\times9$
\item Subgrid: $3\times3$
\item Domain: $\{1..9\}$
\end{itemize}

\section*{Interpretation}
The Sudoku puzzle is embedded into a three dimensional configuration space $(x,y,z)$.
The XY plane preserves board structure while the Z axis allows exploration of alternate configurations.

A valid solution occurs when:
\begin{equation}
\Phi(S)=0
\end{equation}

\end{document}
